\documentclass[tikz]{standalone}
\usepackage{amsmath}
\definecolor{bgcol}{HTML}{171D2B}
\definecolor{accent}{HTML}{F7B538}
\definecolor{setcol}{HTML}{7DA2C9}
\pagecolor{bgcol}
\begin{document}
% Gambar 1: tabung polos — dua elips dan dua garis tegak
% Elips atas utuh, elips bawah: belakang putus-putus, depan penuh
\begin{tikzpicture}[scale=1.0, color=white, font=\large]
  % Elips atas: digambar utuh
  \draw[thick] (0, 2) ellipse (2 and 0.6);
  % Dua garis tegak selimut
  \draw[thick] (-2, 2) -- (-2, -2);
  \draw[thick] (2, 2) -- (2, -2);
  % Elips bawah, bagian belakang (tersembunyi): garis putus-putus
  \draw[thick, dashed] (2, -2) arc (0:180:2 and 0.6);
  % Elips bawah, bagian depan (terlihat): garis penuh
  \draw[thick] (-2, -2) arc (180:360:2 and 0.6);
\end{tikzpicture}

% Gambar 2: tabung berlabel — jari-jari r dan tinggi t
% r = 7 cm, t = 10 cm pada soal penutup artikel
% V = pi r^2 t = 22/7 x 49 x 10 = 1540 cm^3
\begin{tikzpicture}[scale=1.0, color=white, font=\large]
  % Elips atas utuh
  \draw[thick] (0, 2) ellipse (2 and 0.6);
  % Dua garis tegak selimut
  \draw[thick] (-2, 2) -- (-2, -2);
  \draw[thick] (2, 2) -- (2, -2);
  % Elips bawah: belakang putus-putus, depan penuh
  \draw[thick, dashed] (2, -2) arc (0:180:2 and 0.6);
  \draw[thick] (-2, -2) arc (180:360:2 and 0.6);
  % Jari-jari: dari pusat elips atas ke ujung kanan
  \fill[accent] (0, 2) circle (0.06);
  \draw[thick, accent] (0, 2) -- (2, 2);
  \node[accent, above] at (1, 2) {$r$};
  % Tinggi: garis putus-putus di tengah tabung
  \draw[thick, dashed, accent] (0, 2) -- (0, -2);
  \node[accent, right] at (0.15, 0) {$t$};
\end{tikzpicture}
\end{document}
